\section{Related Work}
\label{sec:related-work}

We organize prior work into four streams---silicon sampling and synthetic respondents, persona-based and personalized simulation, LLMs as representation generators, and parameter-efficient adaptation---and position our contribution against each.

\subsection{Silicon Sampling and Synthetic Respondents}
The use of LLMs as synthetic survey respondents has grown rapidly, with early applications demonstrating that persona-conditioned models can reproduce aggregate marginals for subpopulations~\citep{sun2024random,sarstedt2024using}. Subsequent critical work has shown that model answers diverge from human reference distributions in systematic ways: \citet{dominguezolmedo2023questioning} question whether LLM survey responses faithfully reflect demographic distributions, \citet{qu2024performance} document performance and bias gaps in public-opinion simulation, and \citet{taubenfeld2024systematic} expose systematic biases in LLM-simulated debates. Reliability audits of persona-conditioned respondents~\citep{morocho2026assessing} confirm that multi-attribute personas fail to guarantee faithful individual-level prediction. Most relevant to our work, \citet{ku2026silicon} introduce cross-survey transfer evaluation for silicon sampling and show that methods performing well within a single instrument degrade substantially under transfer; our OOD-demographic-intersection regime (Section~\ref{sec:regimes}) directly addresses the transfer-evaluation gap they identify, and our representation-based factorization addresses the structural cause of the degradation they observe. \citet{anthis2025social} and \citet{bail2024generative} survey the broader promise and methodological challenges of LLM-based social simulation, while \citet{westwood2025potential} frames LLMs as a potentially existential threat to online survey research, motivating the need for methods that augment rather than replace human respondents. The AI-Augmented Surveys framework of \citet{kim2023aiaugmented} is the closest task-level precedent: they use LLMs to predict opinions for unasked items, but they do so by direct generation and without the representation--decoder factorization we propose.

\subsection{Persona-Based and Personalized Simulation}
A large literature conditions LLMs on personas to simulate individuals or populations. \citet{hu2024quantifying} quantify the persona effect---the gain from adding demographic, social, and behavioral variables---and inform our demographic-block design. \citet{li2026personabased} scale persona-based simulation to population-level opinion modeling, while \citet{venkit2026need} argue that coarse sociodemographic personas are insufficient and call for socially-grounded persona frameworks; our history block, which retains option-level revealed preferences, is a concrete response to this call. \citet{su2026adaptive} show that evidence-grounded adaptive interviewing improves persona--decision alignment, supporting our use of retrieved answer history as evidence. \citet{li2025generated} caution that LLM-generated personas are ``a promise with a catch,'' with reliability varying sharply by attribute, and \citet{tseng2024tales} survey the dual role of persona in role-playing and personalization. On the personalization-systems side, \citet{sun2024personadb} introduce Persona-DB for efficient LLM personalization in response prediction via collaborative data refinement, which we adopt as a strong baseline (B3). \citet{chen2024when} survey personalization challenges and opportunities for LLMs more broadly. Conceptually closest to our RespondentVec, \citet{beckmann2026where} extract ``persona vectors'' via mechanistic interpretability to individuate LLMs, but their analysis is philosophical and mechanistic rather than predictive; our RespondentVec operationalizes a similar intuition as a reusable predictor. Crucially, all of the above operate in the direct-generation paradigm; none factorizes prediction into a latent representation plus a calibratable decoder, which is our structural departure.

\subsection{LLMs as Representation Generators}
Our central analogy draws on \citet{he2025geolocation}, who show that LLM-derived geolocation representations serve as generic enhancers for spatio-temporal learning, dramatically outperforming direct LLM forecasting by decoupling the LLM's location semantics from the downstream temporal model. We transpose this predictor-to-representation-generator move to the survey domain, treating the respondent's latent reaction rather than the location as the object to be represented. The technical engine for extraction is LLM2Vec-Gen~\citep{behnamghader2026llmvecgen}, a generative embedding method that preserves the LLM's output distribution semantics during contrastive alignment, in contrast to the discriminative LLM2Vec~\citep{behnamghader2024llmvec} which discards them; we exploit this property to ensure the reaction vector carries generative disposition toward answers. Related generative-representation approaches include GRACE~\citep{sun2025grace}, which uses contrastive policy optimization, and generative representational instruction tuning~\citep{muennighoff2024generative}. We also draw on the theoretical limitations of embedding-based retrieval~\citep{weller2025theoretical} to motivate the option-aware decoder, which restores instrument-specific structure that a pure embedding cannot capture. The broader literature on LLM-based embedding models~\citep{lee2024nvembed,wang2024improving,wang2024multilingual,enevoldsen2025mmteb} and latent reasoning via embedding prediction~\citep{hwang2025latent} informs our choice to extract representations rather than answer tokens, but these works target retrieval or reasoning rather than respondent simulation.

\subsection{Parameter-Efficient Adaptation and Generative Agents}
The QLoRA baseline (B5) draws on the large literature on low-rank adaptation~\citep{mao2024survey}, including personalized RLHF with shared adapters~\citep{liu2025shared} and personalized driver-assistance fine-tuning~\citep{song2025enhancing}, which represent the task-adaptivity upper bound at high compute. The generative-agent baseline (B6) draws on the influential Generative Agents framework~\citep{park2023generative} and its general-purpose individual-simulation extension~\citep{park2024agents}, which ground agents in self-reports; related multi-agent and society-level simulations~\citep{gao2024large,hou2025society} provide context for the agent paradigm. None of these methods, however, separates representation from prediction in the way we propose, and our cost--accuracy analysis (Section~\ref{sec:tradeoff}) shows that the per-respondent fine-tuning cost of QLoRA and the per-item memory-retrieval cost of generative agents both sit far from the Pareto frontier that our representation-based methods trace.
